\documentclass[12pt]{article}

\usepackage{fontspec}
\newfontfamily\handwriting{a-casual-handwritten-pen}[Path=../assets/fonts/,Extension=.ttf,Scale=1.5]
\newfontfamily\comic{Comic Neue}

\usepackage[utf8]{inputenc}
\usepackage{amsmath}
\usepackage{amssymb}
\usepackage{enumitem}
\usepackage{xcolor}
\usepackage{soul}
\usepackage{tikz}
\usepackage[most]{tcolorbox}
\usepackage{titlesec}
\usepackage{multicol}
\usepackage{environ}

\usetikzlibrary{fit, positioning, arrows.meta}

% --- Custom Colors ---
\definecolor{primaryblue}{RGB}{42, 111, 151}
\definecolor{exbg}{RGB}{255, 240, 245}
\definecolor{rulebg}{RGB}{232, 245, 233}
\definecolor{highlight}{RGB}{255, 243, 176}
\definecolor{pinkanno}{RGB}{255, 105, 180}
\definecolor{purpleanno}{RGB}{147, 112, 219}
\definecolor{solutionbg}{RGB}{245, 245, 255}

% --- Header Styling ---
\titleformat{\section}{\color{primaryblue}\normalfont\Large\bfseries}{\thesection}{1em}{}
\titleformat{\subsection}{\color{primaryblue}\normalfont\large\bfseries}{\thesubsection}{1em}{}

% --- Friendly Boxes ---
\newtcolorbox{friendlyex}[1]{
    colback=exbg,
    colframe=pinkanno!50,
    arc=5pt,
    title=\textbf{#1},
    coltitle=black,
    attach title to upper,
    after title={:\quad}
}

\newtcolorbox{friendlyrule}[1]{
    colback=rulebg,
    colframe=primaryblue!30,
    arc=5pt,
    fonttitle=\bfseries,
    title={#1},
    coltitle=primaryblue
}

\newcommand{\funhl}[1]{\sethlcolor{highlight}\hl{#1}}

% --- Show/Hide Solutions ---
\newif\ifshowsolutions

% Uncomment the next line to show solutions.
\showsolutionstrue

\newcommand{\solutionblank}{\vspace{25ex}}

\NewEnviron{solutionbox}{
\ifshowsolutions
\begin{tcolorbox}[
    colback=solutionbg,
    colframe=purpleanno!45,
    arc=5pt,
    title={Solution},
    coltitle=primaryblue,
    fonttitle=\bfseries
]
\BODY
\end{tcolorbox}
\else
\solutionblank
\fi
}

\begin{document}
\handwriting

\begin{center}
    \Large\textbf{\color{primaryblue}Module 1 Lesson 2\\Domain and Range} \\
    \vspace{0.2cm}
    \hrule
\end{center}

\section*{Objectives}

\begin{friendlyrule}{In this lesson, we will learn how to}
\begin{itemize}
    \item read and write interval notation;
    \item determine the domain of a function algebraically;
    \item determine the domain and range of a function from its graph.
\end{itemize}
\end{friendlyrule}

\section*{1. Interval Notation}

\funhl{Interval notation} is a shorthand for describing a set of real numbers using inequalities.

\begin{friendlyrule}{Interval Notation}
\begin{center}
\begin{tabular}{lll}
Inequality & Interval notation & \\ \hline
$a<x<b$ & $(a,b)$ & open on both ends \\
$a\le x\le b$ & $[a,b]$ & closed on both ends \\
$a\le x<b$ & $[a,b)$ & closed at $a$, open at $b$ \\
$a<x\le b$ & $(a,b]$ & open at $a$, closed at $b$ \\
$a<x$ & $(a,\infty)$ & \\
$a\le x$ & $[a,\infty)$ & \\
$x<b$ & $(-\infty,b)$ & \\
$x\le b$ & $(-\infty,b]$ &
\end{tabular}
\end{center}

A parenthesis means the endpoint is \emph{not} included; a bracket means the endpoint \emph{is} included. Infinity always gets a parenthesis, since it is not a number that can be ``reached.''
\end{friendlyrule}

\section*{2. Domain of a Function}

\begin{friendlyrule}{Domain}
The \funhl{domain} of a function is the complete set of input values ($x$-values) for which the function is defined and produces a real number output. Common restrictions include:
\begin{itemize}
    \item[(a)] avoiding division by zero in rational functions;
    \item[(b)] avoiding negative values under square roots.
\end{itemize}
\end{friendlyrule}

\begin{friendlyex}{Example 1}
Determine the domain of each function.

(Clue: is there any value that $x$ cannot be?)

\begin{itemize}
    \item[(i)] $f(x)=3x^2+2x-1$
    \item[(ii)] $g(x)=\dfrac{5}{x-3}$
    \item[(iii)] $k(x)=\sqrt{3x-12}$
    \item[(iv)] $h(x)=\dfrac{2x}{x^2-x-6}$
\end{itemize}

Challenging one:

\[
h(x)=\frac{3}{\sqrt{2x-5}}.
\]
\end{friendlyex}

 

\begin{solutionbox}
\textbf{(i)} $f$ is a polynomial, so there are no restrictions on $x$. The domain is

\[
(-\infty,\infty).
\]

\textbf{(ii)} We need $x-3\neq 0$, i.e., $x\neq 3$. The domain is

\[
(-\infty,3)\cup(3,\infty).
\]

\textbf{(iii)} We need $3x-12\ge 0$, i.e., $x\ge 4$. The domain is

\[
[4,\infty).
\]

\textbf{(iv)} Factor the denominator:

\[
x^2-x-6=(x-3)(x+2).
\]

We need $x\neq 3$ and $x\neq -2$. The domain is

\[
(-\infty,-2)\cup(-2,3)\cup(3,\infty).
\]

\textbf{Challenging one:} Since the square root is in the denominator, we need the radicand to be \emph{strictly} positive:

\[
2x-5>0\ \Longrightarrow\ x>\frac52.
\]

The domain is

\[
\left(\frac52,\infty\right).
\]
\end{solutionbox}

  

\section*{3. Range of a Function}

\begin{friendlyrule}{Range}
The \funhl{range} is the complete set of output values ($y$-values) generated by applying the function to the inputs in the domain.
\end{friendlyrule}

\begin{friendlyex}{Example 2}
Determine the domain and range of a function whose graph has a vertical asymptote at $x=3$ and a horizontal asymptote at $y=1$.
\end{friendlyex}

 

\begin{solutionbox}
The graph never touches the vertical line $x=3$, so $3$ is excluded from the domain:

\[
\text{domain}=(-\infty,3)\cup(3,\infty).
\]

The graph gets arbitrarily close to $y=1$ but never reaches it, so $1$ is excluded from the range:

\[
\text{range}=(-\infty,1)\cup(1,\infty).
\]
\end{solutionbox}

  

\begin{friendlyex}{Example 3}
Use the pictured function $f$ to find:

\begin{multicols}{2}
\begin{enumerate}[label=(\alph*)]
    \item $f(2)$
    \item $f(1)$
    \item all $x$ for which $f(x)=2$
    \item all $x$ for which $f(x)=-1$
    \item the $x$-intercept(s)
    \item the $y$-intercept
    \item the domain of $f$
    \item the range of $f$
\end{enumerate}
\end{multicols}

The graph consists of: an open circle at $(-3,4)$ connected by a line segment down to $(-2,2)$; a horizontal segment from $(-2,2)$ to $(1,2)$; and a decreasing line segment starting at $(1,2)$, passing through $(3,0)$, and continuing (with an arrow) indefinitely to the lower right.
\end{friendlyex}

 

\begin{solutionbox}
From $(-3,4)$ (open) to $(-2,2)$, the piece has slope

\[
\frac{2-4}{-2-(-3)}=\frac{-2}{1}=-2,
\]

so this piece is $f(x)=-2x-2$.

From $(1,2)$ onward, the piece has slope

\[
\frac{0-2}{3-1}=\frac{-2}{2}=-1,
\]

so this piece is $f(x)=-x+3$.

\textbf{(a)} $x=2$ is on the decreasing piece: $f(2)=-(2)+3=1$.

\textbf{(b)} $x=1$ is on the flat piece: $f(1)=2$.

\textbf{(c)} $f(x)=2$ for every $x$ on the flat segment: $x\in[-2,1]$.

\textbf{(d)} On the decreasing piece, $-x+3=-1$ gives $x=4$.

\textbf{(e)} The graph crosses the $x$-axis where $-x+3=0$, i.e., at $(3,0)$.

\textbf{(f)} $x=0$ is on the flat piece, so the $y$-intercept is $(0,2)$.

\textbf{(g)} The graph starts (excluding) at $x=-3$ and continues without bound to the right:

\[
\text{domain}=(-3,\infty).
\]

\textbf{(h)} The largest $y$-value approached (but not attained) is $4$; the graph decreases without bound to the right:

\[
\text{range}=(-\infty,4).
\]
\end{solutionbox}

  

\section*{Summary}

\begin{friendlyrule}{Key Ideas}
\begin{itemize}
    \item Interval notation uses parentheses for excluded endpoints and brackets for included endpoints; $\infty$ always gets a parenthesis.
    \item The domain of a function excludes values that cause division by zero or a negative number under a square root.
    \item The range is the set of all output values the function actually produces.
    \item Domain and range can often be read directly from a graph by looking at its horizontal and vertical extent.
\end{itemize}
\end{friendlyrule}

\end{document}
